\documentclass{article}
\usepackage{graphicx} % Required for inserting images
\usepackage{amsmath}
\usepackage{hyperref}
\usepackage{authblk} % For affiliations

\usepackage[a4paper, margin=1in]{geometry}

\title{Exploring AIGS: From System Design to Experimental Trade-offs in AI-Generated Science}
\author[1]{Peng Li}

\affil[1]{Institute for AI Industry Research (AIR), Tsinghua University}
\date{October 2024}

% 整体写作思路：

% 1. Paper 要展示的内容：
%   1.1. AIGS 这个概念，通用的 AIGS System 长什么样，当前的 AIGS System 缺少什么
%   1.2. DSL 抽象的通用性及相关讨论
%   1.3. Ablation Module 的必要性及相关讨论
%   1.4. 执行实验过程中发现的各种 Tradeoff（Multi-Sampling 和 Cost 之间；Idea Novelty 和 Execution 之间；Ablation 中）

% 2. 写作思路：
%   2.1. Abstract
%   2.2. Introduction
%   2.3. How Human Conduct Scientific Research
%       描述人类是如何执行科学研究的，对科学研究进行建模
%   2.4. Literature Review for AIGS
%       对过去的论文进行分类总结
%           - Foundational Perspectives on AI in Scientific Discovery
%           - AI as a Research Assistant: Tools for Human-AI Collaboration
%           - Task-Specific AI for Performance Optimization
%           - Autonomous AI in Scientific Discovery: The Automated Scientist
%   2.5. AIGS System Design
%       描述一个通用的 AIGS System Design（但考虑到实现，忽略了 2.3 中的某几条）
%           - DSL Design / Experiment 结合
%           - Proposal Agent
%           - Review Agent
%           - Ablation Agent
%   2.6. Proof-of-Concept Studies
%       实验证明
%           - RAG
%               - Experiment Setup
%               - Results and Analysis
%           - NanoGPT
%           - Self-Instruct
%   2.7. Discussion and Actionable Insights
%           - The Generalizability of DSL Abstractions
%           - The Necessity of the Ablation Module
%           - Trade-offs Identified During Experimental Execution
%               - Agent Flexibility, Idea Novelty and Execution Feasibility
%               - Multi-Sampling Strategies and Cost Efficiency
%               - Insights from Ablation Trials
%   2.8. Limitations and Ethical Considerations
%           - Incompleteness of System Design
%           - Incompleteness of Task Selection
%           - Detailed Examination of System Modules
%   2.9. Future Plan

\begin{document}

\maketitle

\begin{abstract}

% 思路：写作要素：
%   - 科学知识变化日新月异，科研人员压力增加；
%   - Agent 发展引发人们关注，AI 自主生成科学知识得到希望
%   - 文章的目的：
%       - 对科研场景进行建模；
%       - 提出一种在当前 LLM 能力条件下的通用的 AIGS 框架；
%       - 通过实验探索为未来 AIGS 的发展提出 Insights

\end{abstract}

% 【图表】：一张图展示 AIGS 这一概念（e.g. 半张图展示概念，半张图展示我们的 System）

\section{Introduction}

% 思路：交代清楚当前的背景，介绍文章行文思路（每一个 section 做什么），引出后续内容即可
%   - 背景（LLM，Agent 发展迅速，自主智能体备受关注；然而科研人员压力日渐增大，使用 Agent 自主生成科学发现开始被人们关注）
%   - 介绍我们的行文思路：
%       - 本文将首先介绍人类科学家如何进行科学发现，尝试对科学研究进行建模；
%       - 然后对当前 AIGS 领域的研究进行 Review；
%       - 之后提出一个 AIGS 系统设计方案；
%       - 在三种场景上对我们提出的系统进行验证研究；
%       - 深入讨论我们在系统设计以及实验过程中的发现；
%       - 工作的局限性以及未来工作。

\section{How Human Scientist Conduct Scientific Research}

% 思路：TODO 感觉可能得尽可能讨论的详细一些，但是又能收回到我们自己提出的系统

% 【图表】配合展示

\section{Literature Review for AIGS}

% 思路：介绍我们的分类依据，分类 Metric

% 【图表】一张大表统计过去的文章的类型，Metric：
%   - Copilot/Autopilot
%   - 是否自主提出 proposal
%   - 是否有 review 过程
%   - 是否 iter
%   - 是否有 ablation
%   - 能否提出一些优化的方法
%   - 能否提出一些 science knowledge
%   - 是否对提出的 method/science knowledge 进行进一步验证
%   - 人是否有创造性劳动
%   - ...

\subsection{Foundational Perspectives on AI in Scientific Discovery}

% 思路：介绍当前的基础科研工作（例如 Can LLMs Generate Novel Research Ideas?; Can Large Language Models Unlock Novel Scientific Research Ideas? 等）

\subsection{AI as a Research Assistant: Tools for Human-AI Collaboration}

% 思路：AI 辅助人类科研的一系列工作

\subsection{Task-Specific AI for Performance Optimization}

% 思路：专注于提高任务性能表现的工作

\subsection{Autonomous AI in Scientific Discovery: The Automated Scientist}

% 思路：能够自主完成 pipeline，可以提出一些 proposal 的工作

\section{AIGS System Design}

% 思路：我们的目的：为了验证当前 AIGS 工作往往缺失的模块的必要性，提出一个满足我们理想的 AIGS 的最小 Pipeline（可以自主迭代优化 proposal，自主完成实验并提出反馈，自主提出 ablation study 验证 science knowledge）；
%   - 然后基本概括一下我们这个 System（要强调 History 这一要素），讲一下之后要从哪几个部分开始写

% 【图标】需要一张图来展示我们的 System

\subsection{DSL Design}

% 思路：包含的要素：
%   - DSL 设计的 Motivation；
%   - DSL 的目的：作为 Proposal 与 Experiment 之间的桥梁。

\subsection{Proposal Agent}

% 思路：包含的要素：
%   - Proposal Agent 基本功能；
%   - Multi-Sampling 的设计模式。

\subsection{Review Agent}

% 思路：包含的要素：
%   - Review Agent 基本功能；
%   - Review 的内容（Feature 的设计等）。

\subsection{Ablation Agent}

% 思路：包含的要素：
%   - Ablation Agent 基本功能；
%   - 目的，设计思想？

\section{Proof-of-Concept Studies}

% 思路：介绍为了验证我们的系统，挑选了哪些任务来进行验证；我们每个任务挑选的依据以及其目的是什么。

\subsection{RAG} % 可能得换个名字

\paragraph{Experiment Setup}

\paragraph{Results and Analysis}

\subsection{NanoGPT} % 可能得换个名字

\paragraph{Experiment Setup}

\paragraph{Results and Analysis}

\subsection{Self-Instruct} % 可能得换个名字

\paragraph{Experiment Setup}

\paragraph{Results and Analysis}

\section{Discussion and Actionable Insights}

% 思路：用于整理实验过程中的发现，并期望输出一些有价值的 Insights

\subsection{The Generalizability of DSL Abstractions}

% 思路：DSL 抽象的讨论

\subsection{The Necessity of the Ablation Agent}

% 思路：对 Ablation Agent 的讨论

\subsection{Trade-offs Identified During Experimental Execution}

% 思路：讨论实验过程中发现的各种 Trade-off

\paragraph{Agent Flexibility, Idea Novelty and Execution Feasibility}

% 思路：Agent 自由度，Idea 新颖程度与执行的可行性之间的 Trade-off

\paragraph{Multi-Sampling Strategies and Cost Efficiency}

% 思路：多采样鸟枪法与消耗之间的 tradeoff

\paragraph{Insights from Ablation Trials}

% 思路：Ablation 实验过程中的自由度带来的 Trade-off

\subsection{Challenges and Technical Bottlenecks in end-to-end scientific discovery}

% 一个有价值的 insights 是目前 AIGS 系统中（每个部分）遇到的主要难点和挑战，例如模型写代码能力不足导致支撑不了实验。就目前的研究现状（刚刚起步）而言，这些“失败的经验”也是一块 contribution

\section{Limitations and Ethical Considerations}

% 思路：实验的限制，这包含：
%   - 系统设计不够完善，例如缺少文献等模块；
%   - 任务选择不够完整，并且最终没有很哇塞的科学发现；
%   - 缺少对每个模块的详细分析与深入研究。

\subsection{Incompleteness of System Design}

\subsection{Incompleteness of Task Selection and Final Science Generated}

\subsection{Detailed Examination of System Modules}

\section{Future Plan}

% 思路：介绍未来的计划

\appendix
\section{Appendix}

\end{document}z
